%% TEST-NAME: stress-ta-appendix-gray
%% TEST-PURPOSE: Appendix-mode TA stress fixture with gray backrefs and custom page-number formatting.
%% TEST-KIND: stress
%% TEST-BEHAVIOR: B-NUM-SHARED, B-REND-APPENDIX, B-LINK-CLICKABLE, B-LINK-APPENDIX-BIDIR, B-LINK-NO-ORPHAN, B-USEDIN-READING-ORDER, B-PRINTKIND-DISPLAY-NAME, B-PRINTKIND-DISPLAY-OVERRIDE
%% TEST-STATUS: EXPLORATORY
%% TEST-RENDER-MODES: appendix, gray-backrefs, custom-page-format
%% TEST-PDF-PAGES: 7
%% TEST-PDF-STEXT: Stress fixture appendix gray
%%
%% W05-A2 RED gate: joint-atom rendering for Section 9 (thm:s9-jatom-x, thm:s9-jatom-y)
%%   All assertions below FAIL on current HEAD — A1 serial attribution produces
%%   separate 9.9* and 9.10* entries; no {9.9, 9.10}* token appears in the PDF.
%%   After A2 D1: one joint atom {9.9, 9.10}* in the Dependency Index.
%%
%% Dependency Index contains the joint atom identifier {9.9, 9.10}*
%% TEST-PDF-STEXT: \{9\.9, 9\.10\}\*
%% W05-PRINTKIND-DISPLAY-NAME: display name in appendix index must be "Printkind Stress", not "Printkindstress".
%% TEST-PDF-STEXT: Printkind Stress 9\.12 &#x2014; p\. 6
%% TEST-PDF-STEXT-NOT: Printkindstress

\def\START{}\def\END{}
\documentclass{article}
%% Active non-default overrides:
%%   - line 22: appendix/pagenum-format={##1~\textemdash~p.~##2}
%%              exercises the API; default is {##1~(p.\@~##2)}
%%   - appendix-display={tadec=Theorem TA}: plain label override for decorated env
%%              exercises the override API; default would show the \thmsymbol glyph
%% (Sentinels: lines tagged %%NONDEFAULT-OVERRIDE%% deviate from package defaults.)
\usepackage[T1]{fontenc}
\usepackage{amsmath,amssymb}
\usepackage{mathtools}
\usepackage{graphicx}
\usepackage{tikz}
\usetikzlibrary{cd}
\usepackage{enumitem}

%% Load the real theorem style from the monograph
\usepackage{sty-theorems-ta}  % loads amsthm, thmtools, xparse, fontawesome, xcolor
\usepackage{thm-restate}

%% W05-PRINTKIND-DISPLAY-NAME stress: exercises \textsc{} macro-token display-name
%% capture.  Must be declared BEFORE \codeptrack so the print-kind resolver can
%% read \thmt@original@printkindstress at registration time.
%% Uses numberwithin (not [sibling]) to avoid a thmtools+sibling counter crash (P02 scope gap).
\declaretheorem[name={\textsc{Printkind Stress}},numberwithin=section]{printkindstress}

%% W05-PRINTKIND-DISPLAY-OVERRIDE stress: ta-style decorated env whose appendix
%% label is overridden to the plain string "Theorem TA", stripping the \thmsymbol
%% glyph from the index row.  Must be declared BEFORE \codeptrack.
\makeTheoremEnv{tadec}{Theorem TA}{\thmsymbol}{\thmsymbol}[theorem]

%% codependent in appendix mode (W04: exercises appendix bidir + page numbers)
\usepackage[backrefs=appendix,backref-color=gray]{codependent}
%% W04-P04: page-number format override — uses non-default value to exercise the API
\codepsetup{appendix/pagenum-format={##1~\textemdash~p.~##2}}  %%NONDEFAULT-OVERRIDE%%
%% W05-PRINTKIND-DISPLAY-OVERRIDE: plain label override for the decorated tadec env
\codepsetup{appendix-display={tadec=Theorem TA}}  %%NONDEFAULT-OVERRIDE%%
\codeptrack{theorem,proposition,lemma,corollary,definition,notation,terminology,example,remark,question,printkindstress,tadec}

%% Multi-file compatibility check: load subfiles and standalone.
%% These must not conflict with codependent's hooks.
\usepackage{subfiles}
\usepackage{standalone}

\usepackage[hypertexnames=false]{hyperref}
\usepackage{cleveref}

%% cleveref names for custom envs that thmtools doesn't auto-register
\crefname{notation}{Notation}{Notations}
\crefname{terminology}{Terminology}{Terminologies}
\crefname{question}{Question}{Questions}
\crefname{printkindstress}{Printkind Stress}{Printkind Stresses}
\crefname{tadec}{Theorem TA}{Theorems TA}

%% ---------- Concept-tracked macros ----------
\codepnewcommand{\Funct}[2]{\mathrm{#1}\colon #2}
\codepnewcommand{\catname}[1]{\mathsf{#1}}
\codepnewcommand{\Tr}[0]{\operatorname{Tr}}

%% ---------- Equation numbering by section ----------
\numberwithin{equation}{section}

\title{Stress fixture appendix gray}
\author{Test}
\date{}

\begin{document}
\START
\maketitle

%% ====================================================================
%%  Section 1: Counter sharing under load
%% ====================================================================
\section{Counter sharing under load}
\label{sec:counters}

Let $\catname{C}$ be a monoidal category with morphisms $\Funct{F}{C}$.

\begin{definition}\label{def:enrichment}
A $\catname{V}$-enriched category $\catname{C}$ consists of a set of
objects $\mathrm{Ob}(\catname{C})$ and, for each pair $(A,B)$, a hom-object
$\catname{V}(A,B)$ in $\catname{V}$.
\codeptag{uid}{cat:enriched-category}%
\end{definition}

\begin{theorem}\label{thm:main}
Let $A$ and $B$ be objects in $\catname*{C}$.  The functor
$\Funct*{G}{D}$ is fully faithful.
\end{theorem}

\begin{proof}
By \cref{def:enrichment}, obvious.
\end{proof}

\begin{corollary}\label{cor:equiv}
If $F \dashv G$ is an adjoint equivalence, then
$\catname{C} \simeq \catname{D}$.
\end{corollary}

\begin{remark}\label{rem:day}
Day convolution endows the functor category
$[\catname{C}^{\mathrm{op}}, \catname{Vect}_k]$ with a monoidal structure.
This follows from \cref{thm:main} and \cref{cor:equiv}.
\end{remark}

\begin{proposition}\label{prop:compose}
Monoidal functors compose.  This extends \cref{thm:main}.
\end{proposition}

\begin{lemma}\label{lem:technical}
A technical lemma used in the proof of \cref{prop:compose}.
\end{lemma}

\begin{notation}\label{not:hom}
We write $[A,B]$ for the internal hom in $\catname{V}$
(\cref{def:enrichment}).
\end{notation}

\begin{terminology}\label{term:monoidal}
A \emph{monoidal category} is a category $\catname{C}$ equipped
with a bifunctor $\otimes$ and a unit object $I$.
\end{terminology}

\begin{example}\label{ex:vect}
The category $\catname{Vect}_k$ of finite-dimensional vector spaces
over a field $k$ is monoidal under $\otimes_k$.
See \cref{term:monoidal}.
\end{example}

\begin{question}\label{q:coherence}
Is every monoidal category (\cref{term:monoidal}) monoidally
equivalent to a strict one?
See \cref{thm:main,cor:equiv,prop:compose}.
\end{question}

A closing paragraph referencing \cref{lem:technical} and \cref{rem:day}.
All ten environments above should be numbered 1.1 through 1.10 sequentially.


%% ====================================================================
%%  Section 2: Display math suppression
%% ====================================================================
\section{Display math suppression}
\label{sec:display}

A paragraph before display math.

\begin{theorem}\label{thm:idempotent}
For all endomorphisms $f$ in $\catname{C}$ (\cref{def:enrichment}), we have
\begin{equation}\label{eq:idem}
  f \circ f = f.
\end{equation}
In particular, $f$ is idempotent (\cref{term:monoidal}).
\end{theorem}

\begin{proof}
By \cref{lem:technical}, consider the aligned computation:
\begin{align}
  f \circ f &= f \circ (f \circ f) \\
            &= f.
\end{align}
\end{proof}

A paragraph after the proof.  The equation and align environments
should NOT get margin atom numbers.


%% ====================================================================
%%  Section 3: Cross-section backrefs
%% ====================================================================
\section{Cross-section backrefs}
\label{sec:backrefs}

\begin{theorem}\label{thm:sec3}
This uses \cref{def:enrichment} and \cref{thm:main} from section~1.
Those atoms should show ``Used in 3.1.''
\end{theorem}

\begin{definition}\label{def:gadget}
A \emph{gadget} extends an enriched category (\cref{def:enrichment}).
Answers \cref{q:coherence} partially.
\end{definition}

\begin{corollary}\label{cor:sec3}
Combines \cref{thm:sec3}, \cref{def:gadget}, and \cref{prop:compose}.
\end{corollary}

Cross-section remote backref read: see also \codepbackrefsof{thm:main}
for related results.

A second cross-section remote read: see also \codepbackrefsof{thm:idempotent}
for display-math dependencies.


\subsection*{Track-2 reading-order stress}

\begin{theorem}\label{thm:stress-tracktwo-target}
Stress target theorem for Track-2 ordering.
\end{theorem}

\begin{lemma}\label{lem:stress-tracktwo-before}
A theorem-source citation encountered before the align cites
\cref{thm:stress-tracktwo-target}.
\end{lemma}

\begin{align}
  U_1 &= V_1 \qquad\text{first stress reading-order row}
    \label{eq:stress-tracktwo-first} \\
  U_2 &= V_2 \qquad\text{second stress reading-order row}
    \label{eq:stress-tracktwo-second} \\
  U_3 &= V_3 \qquad\text{third row cites }\cref{thm:stress-tracktwo-target}
    \label{eq:stress-tracktwo-third}
\end{align}

Stress reverse row references: \ref{eq:stress-tracktwo-third},
\ref{eq:stress-tracktwo-first}, then \ref{eq:stress-tracktwo-second}.

\begin{corollary}\label{cor:stress-tracktwo-after}
A theorem-source citation encountered after the align cites
\cref{thm:stress-tracktwo-target}.
\end{corollary}


%% ====================================================================
%%  Section 4: Advanced features
%% ====================================================================
\section{Advanced features}
\label{sec:advanced}

%% --- 4.1: Restatable theorem (tests the restatable guard) ---
\begin{restatable}{theorem}{CoherenceThm}\label{thm:coherence}
Every monoidal category (\cref{term:monoidal}) is monoidally
equivalent to a strict monoidal category.  Moreover, the
equivalence is compatible with the enrichment of
\cref{def:enrichment}.
\end{restatable}

\begin{proof}
Follows from Mac Lane's coherence theorem and \cref{lem:technical}.
\end{proof}

%% --- 4.2: tikz-cd diagram inside a theorem ---
\begin{proposition}\label{prop:square}
The following square of monoidal functors commutes up to
natural isomorphism:
\begin{equation}\label{eq:square}
  \begin{tikzcd}
    \catname{C}
      \arrow[r, "F"]
      \arrow[d, "S"']
    & \catname{D}
      \arrow[d, "T"]
    \\
    \catname{C}_{\mathrm{str}}
      \arrow[r, "F_{\mathrm{str}}"']
    & \catname{D}_{\mathrm{str}}
  \end{tikzcd}
\end{equation}
\end{proposition}

%% --- 4.3: Concept-tracked macros (def-sites and use-sites) ---
\begin{definition}\label{def:trace}
The \emph{categorical trace} $\Tr*$ of an endomorphism
$f \in \catname{C}(A,A)$ is the composite
\[
  \Tr(f) = \varepsilon \circ (f \otimes \mathrm{id}) \circ \eta.
\]
\codeptag{uid}{cat:trace}%
\end{definition}

\begin{lemma}\label{lem:trace-cyclic}
The trace $\Tr$ satisfies the cyclic property: for all
$f\colon A \to B$ and $g\colon B \to A$ in $\catname{C}$,
\begin{equation}\label{eq:cyclic}
  \Tr(g \circ f) = \Tr(f \circ g).
\end{equation}
\end{lemma}

%% --- 4.4: Separated proof with \codepproofof ---
%%
%% The proof of prop:square is given here, in a different section
%% from the statement.  This tests the \codepproofof mechanism.
\begin{proof}[Proof of Proposition~\ref{prop:square}]\codepproofof{prop:square}%
We construct the natural isomorphism componentwise.  For each
object $X$ in $\catname{C}$, the strictification functor
$\Funct{S}{C}$ sends $X$ to $[X]$, and the square commutes
by the universal property of the tensor product
(\cref{def:enrichment}).
\end{proof}

%% --- 4.5: Equation backref tracking (Track 1 — standalone equation) ---
%%
%% A standalone equation (OUTSIDE a theorem) that references atoms.
%% In outer mode (default), equations outside theorems should produce
%% equation-sourced atomrefs with parenthesised equation numbers.
Functoriality of $F$ gives the following, by \cref{thm:main}:
\begin{equation}\label{eq:functorial}
  F(f \circ g) = F(f) \circ F(g)
  \qquad\text{(\cref{cor:equiv})}
\end{equation}

%% --- 4.6: Equation backref tracking (Track 2 — align range) ---
%%
%% An align block (OUTSIDE a theorem) with refs on arrows.
%% Track 2 should produce a range atomref like (4.4--4.5).
The enrichment structure gives:
\begin{align}
  X \otimes Y &\xrightarrow{\text{\cref{def:enrichment}}} \catname{V}(X,Y) \label{eq:enrich-step1} \\
  \catname{V}(X,Y) &\xrightarrow{\text{\cref{thm:main}}} Z \label{eq:enrich-step2}
\end{align}

%% --- 4.7: Equation inside theorem (outer mode — should NOT track eq) ---
%%
%% In outer mode, an equation inside a tracked theorem should fall
%% through to the theorem's atom number, NOT produce an eq-sourced atomref.
\begin{theorem}\label{thm:with-eq}
By \cref{def:enrichment}, we have
\begin{equation}\label{eq:inside-thm}
  f = g \quad\text{by \cref{lem:technical}}.
\end{equation}
\end{theorem}

%% --- 4.8: Multiple cross-references in one atom ---
\begin{theorem}\label{thm:omnibus}
By combining
\cref{def:enrichment,thm:main,cor:equiv,prop:compose,lem:technical,term:monoidal},
we obtain a complete description of the enriched monoidal structure.
This extends the coherence result of \cref{thm:coherence}.
\end{theorem}

%% --- 4.9: Adjacent proof with cross-references (tests starred backrefs) ---
\begin{proof}
Apply \cref{def:enrichment} and \cref{thm:main} to the diagram
of \cref{prop:square}; the result follows from \cref{lem:trace-cyclic}
and $\Tr \circ \Tr = \Tr$.
\end{proof}


%% ====================================================================
%%  Section 5: Edge cases
%% ====================================================================
\section{Edge cases}
\label{sec:edge}

%% --- 5.1: Long multi-paragraph theorem ---
\begin{theorem}\label{thm:long}
Let $\catname{C}$ be a monoidal category with duals
(\cref{thm:coherence}).  The following
conditions are equivalent.

\begin{enumerate}[label=(\roman*)]
  \item Every object $A$ in $\catname{C}$ is dualisable: there exists
    $A^*$ and morphisms $\eta\colon I \to A \otimes A^*$ and
    $\varepsilon\colon A^* \otimes A \to I$ satisfying the zigzag
    identities.
  \item The functor $\Funct{D}{C}$ sending $A$ to $A^*$ is a
    contravariant equivalence.
  \item The trace $\Tr(f)$ is defined for every endomorphism $f$
    and satisfies the cyclic property (\cref{lem:trace-cyclic}).
\end{enumerate}

Moreover, when these conditions hold, the category $\catname{C}$ is
rigid in the sense of Deligne--Milne.  The proof proceeds by verifying
each implication separately.

The implication (i)$\Rightarrow$(ii) is standard: the evaluation and
coevaluation maps furnish the unit and counit of the adjunction
$A \otimes (-) \dashv A^* \otimes (-)$, and the triangle identities
are precisely the zigzag identities.

For (ii)$\Rightarrow$(iii), note that the duality functor $D$ is
exact when $\catname{C}$ is abelian, and the trace of $f$ can be
computed as $\Tr(f) = \varepsilon \circ (f \otimes \mathrm{id}) \circ \eta$
by the definition (\cref{def:trace}).
\end{theorem}

%% --- 5.2: Theorem ending with display math (no trailing text) ---
\begin{proposition}\label{prop:endmath}
The pentagon identity for the associator $\alpha$
(cf.\ \cref{thm:coherence}) reads
\begin{equation}\label{eq:pentagon}
  (\alpha_{A,B,C} \otimes \mathrm{id}_D)
  \circ \alpha_{A,B \otimes C,D}
  \circ (\mathrm{id}_A \otimes \alpha_{B,C,D})
  = \alpha_{A \otimes B,C,D} \circ \alpha_{A,B,C \otimes D}.
\end{equation}
\end{proposition}

%% --- 5.3: Theorem ending with an itemize list ---
\begin{lemma}\label{lem:endlist}
The following are examples of rigid monoidal categories:
\begin{itemize}
  \item $\catname{Vect}_k$ — finite-dimensional vector spaces (\cref{ex:vect}).
  \item $\catname{Rep}(G)$ — representations of a finite group $G$.
  \item $\catname{Hilb}_{\mathrm{fd}}$ — finite-dimensional Hilbert spaces.
\end{itemize}
\end{lemma}

%% --- 5.4: Nested tracked environment (triggers warning) ---
\begin{definition}\label{def:nested}
A \emph{fusion category} is a rigid semisimple monoidal category
with finitely many simple objects.
\begin{example}\label{ex:nested}
The category $\catname{Rep}(G)$ for a finite group $G$ is fusion.
\end{example}
\end{definition}

%% --- 5.5: Empty section (no content — tests no spurious atoms) ---
\section{Placeholder}
\label{sec:empty}

%% --- 5.6: Restated theorem (must not produce duplicate SBL atom) ---
\section{Restated results}
\label{sec:restate}

We collect restated results for the reader's convenience.

\CoherenceThm*

The restatement above must not produce a duplicate sbl atom record.

A final paragraph referencing multiple labels via cleveref:
\Cref{thm:main,prop:compose,lem:technical} establish the core
categorical framework, while
\cref{cor:equiv,def:enrichment,q:coherence} provide the
foundational vocabulary.  See also \autoref{not:hom} and
\cref{thm:long,prop:endmath,lem:endlist}.


%% ====================================================================
%%  Section 7: W04 proof-anchor features (proof-begin anchormap routing)
%% ====================================================================
\section{Proof-anchor features}
\label{sec:w04proof}

%% --- 7.1: Long-distance proof with starred anchor ---
%%
%% \codepproofof*{label} rebinds this proof to thm:long (section~5).
%% Like the plain form, it now routes backlinks through the proof
%% key's \codep@anchormap entry and the shared proof-begin
%% \Hy@raisedlink emission path, so clicks land near the proof
%% heading rather than on the parent theorem anchor.
\begin{proof}[Proof of Theorem~\ref{thm:long}]\codepproofof*{thm:long}%
By verifying each implication in \cref{thm:long} using the rigidity
criterion of \cref{def:trace} and the cyclic property of \cref{lem:trace-cyclic}.
The dualisability condition (i) implies (ii) by naturality, and (iii)
follows from \cref{def:enrichment}.
\end{proof}

%% ====================================================================
%%  Section 9: Auto-detect proof scenarios (W05-A1 RED gate)
%%  Exercises backward-ref body-cite, forward-ref body-cite,
%%  joint-via-single-cref best-effort-serial (both-tracked + mixed-tracking),
%%  and section-only heading with tracked-entity gate.
%%  After P02+P03 these render correctly; on HEAD they expose routing gaps.
%% ====================================================================
\section{Auto-detect proof scenarios}
\label{sec:w05a1}

%% --- 9.1: Backward-ref body-cite ---
%% Proof heading names thm:s9-backward (defined here); body cites def:trace.
%% After P02: def:trace gets a Used-In tag attributed to proof:9.1.
%% NON-ADJACENT: intervening paragraph clears codep@pendingresultid at para/begin.
\begin{theorem}\label{thm:s9-backward}
Every categorical trace (\cref{def:trace}) satisfies the cyclic
property (\cref{lem:trace-cyclic}).
\end{theorem}

We recall the key definitions before proceeding.

\begin{proof}[Proof of \cref{thm:s9-backward}]
By \cref{def:trace} and the cyclic property \cref{lem:trace-cyclic}.
\end{proof}

%% --- Forward-ref body-cite: proof comes BEFORE the theorem ---
%% Proof heading names thm:s9-forward-late (defined later); body cites
%% def:enrichment.  After P02: unresref:<N> allocates at scan time; body
%% cite migrates to proof:9.2 at drain.  No unresref: survives to final cdp.
\begin{proof}[Proof of \cref{thm:s9-forward-late}]
By \cref{def:enrichment} and monoidal coherence (\cref{thm:coherence}).
\end{proof}

%% thm:s9-forward-late is defined AFTER its proof (forward-ref target).
\begin{theorem}\label{thm:s9-forward-late}
The forward-referenced result, combining coherence (\cref{thm:coherence})
with enrichment (\cref{def:enrichment}).
\end{theorem}

%% --- 9.3, 9.4: Joint-via-single-cref BEST-EFFORT-SERIAL case (a) ---
%% Both heading refs are tracked atoms.  After P02 (sub-delta-7): two
%% independent \codep@cdp@proof records (9.3* and 9.4*).  NOT a single
%% joint atom — A2 deferred.
%% NON-ADJACENT: intervening paragraph clears codep@pendingresultid at para/begin.
\begin{proposition}\label{prop:s9-joint-a}
First joint target for best-effort serial proof attribution.
\end{proposition}

\begin{corollary}\label{cor:s9-joint-b}
Second joint target for best-effort serial proof attribution.
Follows from \cref{prop:s9-joint-a}.
\end{corollary}

We give a unified proof of both results.

\begin{proof}[Proof of \cref{prop:s9-joint-a,cor:s9-joint-b}]
Both \cref{prop:s9-joint-a} and \cref{cor:s9-joint-b} follow from
\cref{lem:trace-cyclic} and monoidal coherence.
\end{proof}

%% --- 9.5: Joint-via-single-cref mixed-tracking case (b) ---
%% One tracked atom (lem:s9-mixed) and one section ref (sec:w05a1).
%% After P02+P03: ONE record for 9.5 (tracked target); tracked-entity
%% gate silently no-ops the section ref.
%% NON-ADJACENT: intervening paragraph clears codep@pendingresultid at para/begin.
\begin{lemma}\label{lem:s9-mixed}
Tracked target for the mixed-tracking joint proof scenario.
\end{lemma}

The tracked lemma and the section reference play different roles here.

\begin{proof}[Proof of \cref{lem:s9-mixed,sec:w05a1}]
The tracked lemma follows from \cref{def:trace}; the section ref
is navigational and triggers the tracked-entity gate.
\end{proof}

%% --- Section-only heading: zero records expected ---
%% After P02+P03: tracked-entity gate fires; sec:w05a1 has only
%% \codep@lblnum, no \codep@lblentity; zero attribution, no cdp@proof.
\begin{proof}[Proof of \autoref{sec:w05a1}]
Section-only heading exercises the tracked-entity gate.
Uses \cref{thm:coherence} in the body (body cite dropped by gate).
\end{proof}

%% --- 9.9, 9.10: Joint-atom proof (W05-A2 RED gate) ---
%% After A2: ONE joint atom {9.9, 9.10}* renders in Used-In entries of
%% both thm:s9-jatom-x and thm:s9-jatom-y, and in the Dependency Index.
%% Currently (A1 serial): two independent records 9.9* and 9.10* appear.
%% NON-ADJACENT: no intervening theorem between these two and the proof.
\begin{theorem}\label{thm:s9-jatom-x}
First theorem for joint-atom proof (W05-A2 deferred).
Uses \cref{def:enrichment}.
\end{theorem}

\begin{theorem}\label{thm:s9-jatom-y}
Second theorem for joint-atom proof (W05-A2 deferred).
Follows from \cref{lem:trace-cyclic}.
\end{theorem}

We give a joint proof of both theorems above.

\begin{proof}[Proof of \cref{thm:s9-jatom-x,thm:s9-jatom-y}]
Both follow from \cref{def:enrichment} and \cref{lem:trace-cyclic}
by a unified monoidal argument.
\end{proof}

%% --- W05-PRINTKIND-DISPLAY-NAME: display-name capture stress ---
%% env name is "printkindstress"; display name is \textsc{Printkind Stress}.
%% After the fix: appendix shows "Printkind Stress N.M" not "Printkindstress N.M".
\begin{printkindstress}\label{thm:printkindstress}
A stress theorem for the printkind display-name exercise.
Uses \cref{def:enrichment}.
\end{printkindstress}

The above result (\cref{thm:printkindstress}) forces a Dependency Index row
and exercises the display-name capture path.

%% --- W05-PRINTKIND-DISPLAY-OVERRIDE: ta-style decorated env with plain override ---
%% Body heading retains the \thmsymbol glyph decoration by design (normal ta-style).
%% The Dependency Index entry uses the plain override "Theorem TA" without the glyph.
%% Automated gate: TEST-PDF-APPENDIX-ENTRY-TEXT-ONLY in trinity-test.lvt.
%% Visual gate: inspect pdf-out/stress-ta-appendix-gray-* and confirm the index
%% row shows "Theorem TA N.M" without the decoration glyph.
\begin{tadec}\label{thm:tadec}
A ta-style decorated theorem whose Dependency Index label uses the plain override.
Uses \cref{def:enrichment}.
\end{tadec}

The above result (\cref{thm:tadec}) forces a Dependency Index row
and exercises the appendix-display override path.


%% ====================================================================
%%  Dependency Index (appendix mode — W04 closeout)
%%  Emits bidirectional links: body-side dagger \to appendix entry AND
%%  appendix entry label \to effective body anchor.  Each appendix row
%%  owns its own \Hy@raisedlink destination, and page numbers use the
%%  /codep/appendix/pagenum-format key set above.
%% ====================================================================
\codepappendix

\END
\end{document}
